% =============================================================================
% Cover Letter for Journal of Air Transport Management (JATM)
% =============================================================================
\documentclass[11pt,a4paper]{article}
\usepackage[utf8]{inputenc}
\usepackage{geometry}\geometry{margin=2.5cm}
\usepackage{microtype}
\usepackage{parskip}
\pagestyle{empty}

\begin{document}

\noindent
Dr.\ Jos\'e Cristiano Pereira \\
Departamento de Engenharia e Computa\c{c}\~ao \\
Universidade Cat\'olica de Petr\'opolis (UCP) \\
Petr\'opolis, RJ, Brazil \\
\textit{Email:} josecristiano.pereira@ucp.br

\bigskip

\noindent
[DATE]

\bigskip

\noindent
Editor-in-Chief \\
\textit{Journal of Air Transport Management} \\
Elsevier

\bigskip

\noindent Dear Editor,

We are pleased to submit our revised manuscript entitled
\textit{``A Kirkpatrick-based Framework for Integrating Safety and
Sustainability in Aviation Maintenance: A Case Study on Waste Reduction and
Resource Efficiency''}, for consideration as a Case Study in the
\textit{Journal of Air Transport Management}.

The aviation maintenance, repair and overhaul (MRO) sector sits at the
intersection of two pressures that the journal regularly addresses:
safety-driven regulation under ICAO/FAA/EASA/ANAC oversight, and the rising
demand for credible sustainability disclosures from airlines and their
supply chain. Maintenance errors contribute to between 12\% and 18\% of
commercial aviation accidents and simultaneously generate substantial
material, energy and labour waste. Despite the strategic relevance of
technical training as a high-leverage intervention, empirical evidence
on how to operationalise training evaluation beyond Level~2 of the
Kirkpatrick model in AS9100-certified MRO settings is scarce, and the
explicit translation of training-effectiveness gains into cleaner-production
outcomes has not been documented in a single longitudinal case.

Our study addresses this gap through a twelve-month case study at a
multinational aircraft-engine MRO facility. We operationalise the four
Kirkpatrick levels (Reaction, Learning, Behaviour, Results) within the
existing training process, with each level mapped onto specific AS9100
clauses, SMS pillars and sustainability indicators. The intervention
was associated with a 28\% lower training-failure rate and a 35\% lower
regulatory non-conformance rate, with estimated annual quality-cost
savings of USD~0.75--1.5~M across pessimistic-to-optimistic scenarios
(USD~1.19~M base case) and a payback period under twelve months in every
scenario. Customer satisfaction surveys with four major airline customers reported
improvements across all five evaluated categories during the intervention
year. The single-site
case-study design supports association rather than causal attribution,
and we report the results accordingly. We derive managerial implications
for MRO operators and policy recommendations for civil aviation
authorities and industry bodies (IATA, SAE).

We believe this work aligns directly with the editorial scope of JATM in
three respects:

\begin{itemize}
  \item It contributes to the journal's recurring theme of aviation safety
        management, providing measurable evidence of how training
        evaluation translates into regulatory compliance.
  \item It contributes to the growing JATM stream on sustainability in
        air transport, demonstrating an auditable connection between human
        capital interventions and cleaner production.
  \item It offers actionable policy guidance addressed to civil aviation
        authorities and industry bodies, in line with the journal's
        managerial and policy orientation.
\end{itemize}

In response to the previous round of constructive feedback, the revised
manuscript adds a dedicated \textit{Managerial Implications} section, a
\textit{Policy Recommendations} section, an elaborated \textit{Economic
Analysis} of cost-of-quality decomposition, a comparative
\textit{Before/After} results table, a structured \textit{Limitations}
table, and a new conceptual figure linking Kirkpatrick levels to AS9100
clauses and sustainability outcomes. The abstract has been rewritten to
lead with economic outcomes, and the literature review has been
strengthened with recent JATM references.

We confirm that:

\begin{itemize}
  \item This manuscript is original, has not been published previously,
        and is not under consideration for publication elsewhere.
  \item All authors have read and approved the submission.
  \item There are no known conflicts of interest associated with this work.
  \item Aggregated operational data are reported within the manuscript;
        individual-level data are not publicly available due to commercial
        confidentiality.
\end{itemize}

Should you require any additional information, please do not hesitate to
contact me.

\bigskip

\noindent Sincerely,

\bigskip\bigskip

\noindent
Dr.\ Jos\'e Cristiano Pereira \\
\textit{Corresponding author}\\
on behalf of the authors

\end{document}
